% !TEX encoding = UTF-8 Unicode
\documentclass[a4paper, german, 12pt, onecolumn, oneside,bibliography=totoc,listof=totoc]{article}

\usepackage[utf8]{inputenc}
\usepackage[T1]{fontenc}
\usepackage[pdftex]{graphicx} % Einbindung von Grafiken (pdf, png, jpg)
\usepackage{float}            % bietet Option [H] für bombenfestes Verankern
\usepackage{wrapfig}
\usepackage[ngerman]{babel}   % Silbentrennung nach der neuen deutschen Rechtschreibung, z.B.: Sys-tem
%\usepackage{amstext}         % für Klartext via \text{} in Formeln
\usepackage{color}						% http://en.wikibooks.org/wiki/LaTeX/Colors
\usepackage[hang,small,bf]{caption} % http://radio.wihome.net/wiki/?article=LaTeX:captions
\usepackage{amsmath}
\usepackage[safe]{tipa}
\usepackage{subfig}

\usepackage[pdftex]{hyperref}	% für "klickbare" Verlinkung im Dokument
\hypersetup{
    %bookmarks=true,        % show bookmarks bar?
    unicode=false,          % non-Latin characters in Acrobat’s bookmarks
    pdftoolbar=true,        % show Acrobat’s toolbar?
    pdfmenubar=true,        % show Acrobat’s menu?
    pdffitwindow=false,     % window fit to page when opened
    pdfstartview={FitH},    % fits the width of the page to the window
    pdftitle={Diplomarbeit},    % title
    pdfauthor={Tolleiv Nietsch},     % author
    pdfsubject={},   % subject of the document
    pdfcreator={Tolleiv Nietsch},   % creator of the document
    pdfproducer={TU Chemnitz}, % producer of the document
    pdfkeywords={recommender,collective intelligence}, % list of keywords
    pdfnewwindow=true,      % links in new window
    colorlinks=true,       % false: boxed links; true: colored links
    linkcolor=black,          % color of internal links
    citecolor=black,        % color of links to bibliography
    filecolor=black,      % color of file links
    urlcolor=black           % color of external links
}

\usepackage[square]{natbib}
%\bibpunct{(}{)}{;}{a}{,}{,}
\bibliographystyle{natdin}
\usepackage{multicol}

\usepackage[printonlyused]{acronym}
\renewcommand{\bflabel}[1]{\normalfont{\normalsize{#1}}\hfill}

\usepackage{parskip}          % zw. Absätzen: eine knappe Leerzeile statt hängender Einzüge
\usepackage{makeidx}          % Package zur Indexerstellung
\usepackage{todonotes}				% bringt den todo befehl mit
\let\oldtodo = \todo
\renewcommand{\todo}[1]{{\oldtodo[size=\scriptsize]{#1}}}
\def\argmax{\operatorname*{arg\,max}}
\def\argmin{\operatorname*{arg\,min}}

\newenvironment{packed_enumerate}{
\begin{enumerate}
  \setlength{\itemsep}{1pt}
  \setlength{\parskip}{0pt}
  \setlength{\parsep}{0pt}}{\end{enumerate}
}

\renewcommand{\figurename}{Abb.}
%%%%%%%%%%%%%%%%%%%%%%%%%%%%%%%%%%%%%%%%%%%%%%%%%%%%%%%%%%%%%%%%%%%%%%%%%%%%%%
%
% Größenanpassungen
%
\setlength{\unitlength}{1cm}
\setlength{\oddsidemargin}{0cm}
\setlength{\evensidemargin}{0cm}
\setlength{\textwidth}{16.1cm}
\setlength{\topmargin}{-1.2cm}
\setlength{\textheight}{23cm}
\columnsep 0.5cm
%
%%%%%%%%%%%%%%%%%%%%%%%%%%%%%%%%%%%%%%%%%%%%%%%%%%%%%%%%%%%%%%%%%%%%%%%%%%%%%%

\newcommand{\nosubsubsection}[1] {{\textit{#1}}\newline}
\newcommand{\footnoteremember}[2]{
  \footnote{#2}
  \newcounter{#1}
  \setcounter{#1}{\value{footnote}}
} \newcommand{\footnoterecall}[1]{
  \footnotemark[\value{#1}]
} 

\sloppy                       % großzügiger Zeilenumbruch 
\makeindex 

\hyphenation{Grund-elemente}	% definierte Silbentrennung

\begin{document}
\pagenumbering{roman}
\pagestyle{empty}
\begin{titlepage}
\begin{figure}
  \begin{center}
    \hbox to \hsize{%
      \begin{tabular}[m]{c}
        \includegraphics[width=2.5cm]{Abbildungen/tuc_medieninfo_logo.jpg}
      \end{tabular}
      \hfill%
      \begin{tabular}[m]{c}
        Technische Universität Chemnitz \\
        Fakultät für Informatik \\
        Professur Künstliche Intelligenz \\
      \end{tabular}%
    }
  \end{center}
\end{figure}

\begin{center}
\rule{0pt}{0pt}
\vfill
\vfill
\includegraphics[width=8cm]{Abbildungen/TUC_deutsch_SW.pdf}
\vfill
\vfill
\vfill
\vfill

\begin{Huge}
Diplomarbeit \\
\end{Huge}
im Studiengang Angewandte Informatik \\
~\\
\begin{footnotesize}
Vorgelegt von \\
\end{footnotesize}
\begin{large}
Tolleiv Nietsch \\
\end{large}
\begin{footnotesize}
--- \\
\end{footnotesize}

\vfill
\vfill

\begin{large}Skalierbare Item Recommendation in Big-Data und Suchindexen\end{large} \\

%~ \\
%\hrule
%\begin{Large}Subtitle\end{Large}

\vfill

\end{center}
\end{titlepage}

\newpage
~
\newpage

~
\vspace{5cm}


\begin{large}
\textbf{Selbstständigkeitserklärung} \\ \\
\end{large}
Hiermit erkläre ich, dass ich die vorliegende schriftliche Arbeit selbstständig und ohne Benutzung anderer als der angegebenen Quellen und Hilfsmittel angefertigt habe. \\
Die vorliegende Arbeit ist frei von Plagiaten. Alle Ausführungen, die wörtlich oder inhaltlich aus anderen Schriften entnommen sind, habe ich als solche kenntlich gemacht. \\
Diese Arbeit wurde in gleicher oder ähnlicher Form noch bei keinem anderen Prüfer als Prüfungsleistung eingereicht und ist auch noch nicht veröffentlicht.\\

\vfill
\begin{tabular*}{0.5\textwidth}{@{\extracolsep{\fill}}ll}
Vorname: & Tolleiv \\
Name: & Nietsch \\
Matrikelnummer: & 172314 \\
\end{tabular*}

\vfill
\vfill

Chemnitz, den 29.03.2012 \\
\medskip
\medskip

\underline{~~~~~~~~~~~~~~~~~~~~~~~~~~~~~~~~~~~~~~~~}\\
Tolleiv Nietsch\\


\vfill
\begin{large}
\textbf{Betreuung und Prüfung durch:} \\ \\
\end{large}
\begin{tabular*}{\textwidth}{@{\extracolsep{\fill}}ll}
Prof. Dr. Fred Hamker, & Professur Künstliche Intelligenz, TU-Chemnitz \\
Dr. Johannes Steinmüller & Professur Künstliche Intelligenz, TU-Chemnitz \\ 
\end{tabular*}

\begin{footnotesize}
Technische Universität Chemnitz, Fakultät für Informatik \\
Straße der Nationen 62, 09107 Chemnitz
\end{footnotesize}

\newpage

					% und eidesstattliche Erklärung


%\pagestyle{plain} 
~
\vspace{5cm}

\begin{large}
\textbf{Kurzfassung} \\ \\
\end{large}
tbw

\newpage
\setcounter{page}{1}					% Seitenzahl zurück	
\tableofcontents					% Inhaltsverz.
\newpage

\pagenumbering{arabic}
\renewcommand{\baselinestretch}{1.50}\normalsize
\setcounter{page}{1}					% Seitenzahl zurück	

\section{Einleitung}


Hinweis das ``Elemente'' und ``Item'' oft als Synonym verwendet werden, ähnlich wie auch ``Rating'' und ``Bewertung'' gleichzusetzen sind.

Goldberg92 -> Ursprung des Ausdrucks

Ggf. direkt Terminologie ergänzen - gutes Beispiel Goldberg01 (Kap 3.1)

QVC Italia als Usecase vorstellen und einordnen

\todo[color=red]{Hindernis}
\todo[color=orange]{Wichtige Aufgabe}
\todo[color=yellow]{Aufgabe abhängig von anderen}
\todo[color=green]{Mögliche Ergänzung / weitere Anregung}
\todo[color=white]{}
\section{Grundlagen}

	Auf welchen Themen und Techniken baut die Arbeit auf.

\section{Entwurf}

	8-10 Seiten

\section{Realisation}

	8-10 Seiten
	
	Verfeinerung des Entwurfs, Schilderung bei der Umsetzung aus dem vorangegangenen Kapitel

\section{Evaluation}

	Wie wird gemessen, welche Ergebnisse waren zu erwarten, was wurde erreicht. Warum gibt es Abweichungen, welche Probleme enthält die Messmethode.

	\begin{itemize}
	\item Successful session \citep{hb_18,Smyth05alive-user}
	\item Precision /Recall /F1, \citep{rs}[Kap. 7]
	\item User- und Itemcoverage \citep{rs}[S. 183]
	\end{itemize}

\section{Zusammenfassung}

	Abriss der Arbeit, was wurde erreicht bzw. gelernt. An welcher Stellen kann weitergearbeitet werden. 

\newpage
%
\section{Zusammenfassung}\label{sec:results}

%\textit{Abriss der Arbeit, was wurde erreicht bzw. gelernt. An welcher Stellen kann weitergearbeitet werden.}

\subsection{Fazit}

In dieser Arbeit wurden die Möglichkeiten zur Integration von Suchindexen und Empfehlungsdiensten untersucht. Unter Berücksichtigung der Skalierbarkeit wurden zu diesem Zweck zwei mögliche Lösungswege betrachtet. Zum einen wurden die Ergebnisse eines Empfehlungsdienstes genutzt um Suchergebnisse zu personalisieren. Im zweiten Ansatz wurden Möglichkeiten zur direkten Empfehlungsbildung im Suchindex untersucht um beide Dienste vollständig integriert nutzen zu können.

In der, im Rahmen der Arbeit entwickelten, Beispielapplikation wurden beide Personalisierungslösungen hinsichtlich ihrer Qualität und Leistungsfähigkeit evaluiert. Dabei hat sich gezeigt, dass die bekannten und oft untersuchten Algorithmen des kollaborativen Filterns (vgl. Abschnitt \ref{sec:neighborhoods}) bezüglich der Qualität bessere Ergebnisse liefern als die auf Matrixfaktorisierung basierende Personalisierungslösung. Durch die direktere Integration in Apache Solr lieferen diese allerdings erheblich bessere Leistungswerte.

Umgesetzt wurde die Beispielapplikation auf Basis der quelloffenen Software Apache Solr und Apache Mahout. 
\newpage
\subsection{Ausblick}

Obwohl im Rahmen dieser Arbeit zwei Möglichkeiten zur Personalisierung von Suchergebnissen erfolgreich implementiert wurden, verbleiben einige Fragen zur weiteren Betrachtung. Die im Rahmen der Evaluation festgestellten Skalierungsprobleme von Apache Solr bei niedrigem Paralleliersungsgraden, mögliche Probleme der Zwischenspeichereffizienz bei der Personalisierung und die Optimierung der Parameter bei der Modellberechnung bedürfen weiteren Untersuchungen vor einem praktischen Einsatz. Bei der Verwendung der Personalisierung mittels Webservice gilt es das Problem der disjunkte Kanidatenlisten weiter zu untersuchen.

Im Umgang mit dem Besucher bzw. Kunden einer Webseite sind im Zusammenhang mit den vorgestellten Personalisierungskonzepten Aspekte zu möglichen Einflüssen auf die Verkaufsdiversität (vgl. \citep{Fleder09}) und Möglichkeiten zur Integration von Kritikmechanismen (vgl. \citep{hb_13}). umbetrachtet geblieben.

Bei den eingesetzten Algorithmen ergeben sich ebenfalls Anknüpfungspunkte. Im Bereich der Matrixfaktoriersung existieren zahlreiche Erweiterungen zur Integration von implizitem und explizitem Feedback \citep{Joachims05} und zur Beachtung des temporalen Kontextes \citep{Boughareb11}. Daneben stellt die Kombination mehrerer Empfehlungsalgorithmen in einer Lösung, wie sie zum Beispiel bei der Netflix-Price-Competition notwendig war \citep{netflix2012_2}, einen weiteren wichtigen Aspekt zur Steigerung der Empfehlungsqualität dar. Wie \citep{Forbes11} zeigt, gilt dies ebenso  für die Kombination verschiedener Empfehlungskonzepte. 

Mit diesen zahlreichen Aspekten sollte die Integration eines Empfehlungssystems sicherlich als stetiger Prozess betrachtete werden, der durch neue Erkenntnisse und Algorithmen aus der Forschung vorangetrieben werden kann.



%\begin{itemize}
%\item Ggf. ``Lesezeit'' als Maß einbringen bzw. explizite und implizites Feedback (siehe Joachims05 - Abschnitt 2)
%\item Annahme das sich die Präferenzen der Nutzer über die Zeit nicht änder ist ggf. falsch <-> Temporaler Kontext aus Boughareb11 
%\item Auswirkungen auf die Verkaufsdiversität Fleder09
%\item Ausblick ``Rollout'' \citep{netflix2012_2}
%\item Ausblick ``Similarity Search in High Dimensions via Hashing'' (LSH) als mögliche Erweiterung
%\item Ausblick ``Content boosted matrix factorization'' with in Forbes11
%\item Ausblick andere Suchlösungen - Elastik Search etc..
%\end{itemize}


\newpage
\phantomsection
\addcontentsline{toc}{section}{Abkürzungsverzeichnis}
\section*{Abkürzungsverzeichnis}
\begin{acronym}[XXXXXXX]
	\setlength{\itemsep --- }{-\parsep}
  \setlength{\itemsep}{1pt}
  \setlength{\parskip}{0pt}
  \setlength{\parsep}{0pt}
	\acro{CF}{Collaborative Filtering}
	\acro{GFS} {Google Filesystem}
	\acro{MAE}{Mean Average Error}
	\acro{RMSE}{Root Mean Squared Error}
	\acro{IR}{Information Retrieval}
	\acro{SVD}{Singular Value Decomposition}
\end{acronym}

\phantomsection
\addcontentsline{toc}{section}{Abbildungsverzeichnis}
\listoffigures

\interlinepenalty=10000
%\nocite{*}									% auch die nicht verwendeten bibtex-Einträge einblenden
\cleardoublepage						% damit die TOC auf die richtige seite verweist
\phantomsection							% für hyperref um auf die richtige section zu verlinken
\addcontentsline{toc}{section}{Literatur}
%\renewcommand\refname{I Literatur}

{
\renewcommand{\baselinestretch}{1.2}\footnotesize
\bibliography{Literatur}}

\interlinepenalty=100
\appendix
%\newpage
\section{Irgendein Anhang}
\label{app:anhang}


\end{document}
